%% ----------------------------------------------------------------------------
%%  §5.2  Definizione di derivata
%% ----------------------------------------------------------------------------
\section{Definizione di derivata}
\label{sec:def-derivata}

\begin{definizione}[Derivata]\label{def:derivata}\index{derivata}\index{funzione!derivabile}
Sia $f$ definita nell'intervallo aperto $(a,b)$ e sia $x$ un punto di $(a,b)$. Si dice
che la funzione $f$ è \emph{derivabile} nel punto $x$ se \emph{esiste finito il limite
del rapporto incrementale}
\begin{equation}\label{eq:def-derivata}
  \lim_{h\to0}\frac{f(x+h)-f(x)}{h}.
\end{equation}
Tale limite è la \emph{derivata} di $f$ nel punto $x$, e si indica con una delle
seguenti notazioni, fra loro equivalenti:
\begin{equation}\label{eq:notazioni-derivata}
  f'(x), \qquad \frac{df}{dx}(x), \qquad Df(x), \qquad y', \qquad \frac{dy}{dx}, \qquad Dy .
\end{equation}
Si dice che $f$ è \emph{derivabile nell'intervallo aperto} $(a,b)$ se è derivabile in ogni
punto $x\in(a,b)$.
\end{definizione}

Qualche precisazione sulla definizione. Il rapporto incrementale è una funzione della
sola variabile $h$, definita per i valori $h\neq0$ tali che $x+h\in(a,b)$; poiché $x$ è
interno ad $(a,b)$, questo insieme contiene un intorno di $0$ privato di $0$, e il limite
per $h\to0$ ha senso (\cref{def:limite-funzione}). Le notazioni di Leibniz
$\frac{df}{dx}$, $\frac{dy}{dx}$ ricordano che la derivata è il limite di un rapporto tra
l'incremento $\Delta y=f(x+h)-f(x)$ della variabile dipendente e l'incremento
$\Delta x=h$ della variabile indipendente; non sono, per ora, quozienti di numeri. Con il
cambiamento di variabile $\xi=x+h$ si ottiene la forma equivalente, spesso comoda,
\begin{equation}\label{eq:derivata-xi}
  f'(x)=\lim_{\xi\to x}\frac{f(\xi)-f(x)}{\xi-x}
\end{equation}
(le due forme sono equivalenti per il \cref{teo:limite-composta}, applicato con la
funzione interna $g(\xi)=\xi-x$, diversa da $0$ per $\xi\neq x$, e viceversa con
$g(h)=x+h$, diversa da $x$ per $h\neq0$).
Infine, la derivata è una nozione \emph{locale}: se due funzioni coincidono in un intorno
di $x$, i loro rapporti incrementali coincidono per $h$ piccolo, e per la località del
limite (\cref{lem:localita}) esse hanno la stessa derivata in $x$, oppure nessuna delle
due è derivabile.

In alcuni casi è utile considerare, al posto del limite completo per $h\to0$ nella
\eqref{eq:def-derivata}, soltanto il limite destro per $h\to0^{+}$, oppure il limite per
$h\to0^{-}$. Nel primo caso si parla di \emph{derivata destra}, indicata con
$f'_{+}(x)$, nel secondo caso di \emph{derivata sinistra}, indicata con $f'_{-}(x)$. Per
la \cref{prop:dx-sx}, \emph{$f$ è derivabile in un punto interno $x$ se e solo se
esistono finite e uguali la derivata destra e la derivata sinistra}, e in tal caso
$f'(x)=f'_{+}(x)=f'_{-}(x)$. Se $f$ è definita in $[a,b]$, si dice che $f$ è
\emph{derivabile nell'intervallo chiuso} $[a,b]$ se è derivabile in ogni punto $x\in(a,b)$
e inoltre se $f$ ammette derivata destra (finita) nel punto $x=a$ e derivata sinistra (finita) nel punto
$x=b$.

\subsection{I primi esempi}

\begin{esempio}[Funzioni costanti e lineari]\label{es:derivata-lineare}
Iniziamo dalla funzione costante $f(x)=q$, per ogni $x\in\R$, e proviamo che tale
funzione è derivabile su tutto $\R$ e che la derivata è identicamente nulla; infatti il
rapporto incrementale vale costantemente zero, qualunque sia l'incremento $h\neq0$:
\begin{equation}\label{eq:derivata-costante}
  \frac{f(x+h)-f(x)}{h}=\frac{q-q}{h}=\frac0h=0,
\end{equation}
e quindi anche il limite del rapporto incrementale, per $h\to0$, vale zero (il lettore non
cada nell'errore di considerare il limite per $h\to0$ della \eqref{eq:derivata-costante}
una forma indeterminata $\frac00$: il rapporto è la funzione costante $0$, e il suo limite
è $0$). Più in generale verifichiamo che la derivata della funzione $f(x)=mx+q$, con $m$
e $q$ costanti (il cui grafico è una retta), è identicamente uguale a $m$; infatti il
rapporto incrementale vale costantemente $m$, qualunque sia $h\neq0$:
\begin{equation}\label{eq:derivata-retta}
  \frac{f(x+h)-f(x)}{h}=\frac{[m(x+h)+q]-[mx+q]}{h}=\frac{mh}{h}=m .
\end{equation}
Abbiamo già calcolato nel paragrafo precedente la derivata della funzione $f(x)=x^{2}$,
trovando $f'(x)=2x$ \eqref{eq:tasso-t2}.
\end{esempio}

\begin{esempio}[Il valore assoluto non è derivabile in $0$]\label{es:modulo-non-derivabile}
Verifichiamo invece che la funzione $f(x)=\abs x$ non è derivabile per $x=0$. Infatti, se
$h\neq0$, si ha
\begin{equation}\label{eq:rapporto-modulo}
  \frac{f(0+h)-f(0)}{h}=\frac{\abs{0+h}-\abs0}{h}=\frac{\abs h}{h}.
\end{equation}
Abbiamo già incontrato questa funzione nella \eqref{eq:segno}. Il limite per $h\to0$ del
rapporto incrementale non esiste, perché risulta \eqref{eq:modulo-su-x}
\begin{equation}\label{eq:derivate-modulo}
  \lim_{h\to0^{+}}\frac{\abs h}{h}=1, \qquad \lim_{h\to0^{-}}\frac{\abs h}{h}=-1 .
\end{equation}
Quindi $f(x)=\abs x$ non è derivabile per $x=0$, mentre esistono le derivate destra e
sinistra, uguali rispettivamente a $+1$ e $-1$. Nei punti $x\neq0$ la funzione coincide
localmente con $x$ o con $-x$, ed è quindi derivabile con derivata $1$ per $x>0$ e $-1$
per $x<0$. Geometricamente, il grafico ha in $0$ uno \emph{spigolo}: le due semirette
hanno pendenze diverse (\cref{fig:non-derivabili}, a sinistra).
\end{esempio}

Un altro modo in cui la derivabilità può venire meno è che il rapporto incrementale
tenda all'infinito. Se il limite \eqref{eq:def-derivata} vale $+\infty$ o $-\infty$, si
dice talvolta che $f$ ha in $x$ \emph{derivata infinita}; ma la funzione \emph{non} è
derivabile, perché la definizione richiede un limite finito.

\begin{esempio}[Derivate infinite]\label{es:derivate-infinite}
Per $f(x)=\sqrt[3]x$ (la radice cubica, definita su tutto $\R$ come inversa della
funzione $x^{3}$) si ha, per $h\neq0$,
\[
  \frac{\sqrt[3]h-0}{h}=\frac1{\sqrt[3]{h^{2}}}\to+\infty \qquad (h\to0),
\]
e il grafico ha nell'origine una tangente verticale. Per $f(x)=\sqrt{\abs x}$ si ha invece
$\frac{\sqrt{\abs h}}{h}\to+\infty$ per $h\to0^{+}$ e $\to-\infty$ per $h\to0^{-}$: il
grafico presenta una \emph{cuspide} (\cref{fig:non-derivabili}). In entrambi i casi la
funzione è continua in $0$ ma non derivabile.
\end{esempio}

\begin{figure}[htb]
\centering
\begin{tikzpicture}
\begin{axis}[name=a1, width=0.36\textwidth, height=0.32\textwidth, axis lines=middle,
  xmin=-1.3, xmax=1.3, ymin=-1.2, ymax=1.3, ticks=none, title={\scriptsize $\abs x$: spigolo}]
  \addplot[curva, domain=-1.2:1.2, samples=3] {abs(x)};
\end{axis}
\begin{axis}[at={(a1.east)}, anchor=west, xshift=0.6cm, name=a2, width=0.36\textwidth,
  height=0.32\textwidth, axis lines=middle, xmin=-1.3, xmax=1.3, ymin=-1.2, ymax=1.3,
  ticks=none, title={\scriptsize $\sqrt[3]x$: tangente verticale}]
  \addplot[curva, domain=0:1.2, samples=80] {x^(1/3)};
  \addplot[curva, domain=-1.2:0, samples=80] {-(-x)^(1/3)};
\end{axis}
\begin{axis}[at={(a2.east)}, anchor=west, xshift=0.6cm, width=0.36\textwidth,
  height=0.32\textwidth, axis lines=middle, xmin=-1.3, xmax=1.3, ymin=-1.2, ymax=1.3,
  ticks=none, title={\scriptsize $\sqrt{\abs x}$: cuspide}]
  \addplot[curva, domain=-1.2:1.2, samples=161] {sqrt(abs(x))};
\end{axis}
\end{tikzpicture}
\caption{Tre funzioni continue ma non derivabili in $0$: derivate destra e sinistra
diverse (spigolo), rapporto incrementale che tende a $+\infty$ (tangente verticale),
rapporto incrementale che tende a $+\infty$ da destra e a $-\infty$ da sinistra
(cuspide).}
\label{fig:non-derivabili}
\end{figure}

\subsection{Derivabilità e continuità}

Confrontiamo la nozione di derivabilità con quella di continuità. Ricordiamo che una
funzione $f$ è continua in un punto $x$ se (riprendiamo la \eqref{eq:continua-h})
\begin{equation}\label{eq:continua-h-bis}
  \lim_{h\to0}f(x+h)=f(x).
\end{equation}
Gli esempi precedenti mostrano che una funzione continua può non essere derivabile.
Invece vale il viceversa.

\begin{teorema}[La derivabilità implica la continuità]\label{teo:derivabile-continua}
Ogni funzione derivabile in un punto $x$ è continua in $x$.
\end{teorema}

\begin{schema}
\Ipotesi{$f$ derivabile in $x$, cioè esiste finito $f'(x)$.}
\Tesi{$f(x+h)\to f(x)$ per $h\to0$.}
\end{schema}

\begin{strategia}
L'incremento $f(x+h)-f(x)$ si scrive come prodotto del rapporto incrementale, che tende
al numero finito $f'(x)$, per l'incremento $h$, che tende a zero: il prodotto tende a
zero. Il punto essenziale è la \emph{finitezza} della derivata: se il rapporto
incrementale tendesse all'infinito, avremmo una forma indeterminata $\infty\cdot0$.
\end{strategia}

\begin{dimostrazione}
Per $h\neq0$ si ha l'identità
$f(x+h)=f(x)+\frac{f(x+h)-f(x)}{h}\cdot h$; quindi, per le operazioni con i limiti
(\cref{teo:operazioni-funzioni}),
\begin{equation}\label{eq:derivabile-continua}
\begin{split}
  \lim_{h\to0}f(x+h)&=f(x)+\lim_{h\to0}\bigl[f(x+h)-f(x)\bigr]\\
  &=f(x)+\lim_{h\to0}\frac{f(x+h)-f(x)}{h}\cdot\lim_{h\to0}h\\
  &=f(x)+f'(x)\cdot0=f(x). \qedhere
\end{split}
\end{equation}
\end{dimostrazione}

Il teorema non si inverte, come mostrano le funzioni della \cref{fig:non-derivabili};
anzi, la funzione $\sqrt[3]x$ mostra che un rapporto incrementale infinito è compatibile
con la continuità, e la funzione $\sqrt{\abs x}$ che lo sono anche due rapporti
incrementali unilaterali infiniti di segno opposto. In forma contronominale il teorema
fornisce un criterio utile: \emph{una funzione discontinua in un punto non è derivabile
in quel punto}.

\begin{notastorica}[Una funzione continua senza derivata in alcun punto]
Fino a metà Ottocento molti matematici ritenevano che una funzione continua dovesse
essere derivabile tranne al più in punti isolati, come lo spigolo di $\abs x$.
Nel 1872 Karl Weierstrass presentò all'Accademia delle Scienze di Berlino la funzione
\[
  W(x)=\sum_{n=0}^{\infty}a^{n}\cos(b^{n}\pi x), \qquad 0<a<1,\ b \text{ intero dispari},\ ab>1+\tfrac32\pi,
\]
continua su tutto $\R$ ma \emph{priva di derivata in ogni punto}: il suo grafico ha
spigoli a ogni scala. L'esempio, che Bolzano aveva in sostanza anticipato in lavori
rimasti inediti, sconcertò i contemporanei (Hermite parlò di «piaga lamentevole delle
funzioni continue senza derivata»), ma mostrò definitivamente che continuità e
derivabilità sono concetti ben distinti. Un secolo dopo grafici di questo tipo sarebbero
diventati familiari come \emph{frattali}.
\end{notastorica}

\subsection{Il calcolo numerico della derivata}

La definizione suggerisce un modo immediato per calcolare una derivata al computer:
scegliere un $h$ «piccolo» e calcolare il rapporto incrementale. Quanto piccolo? Il
programma che segue lo fa per $f(x)=\sen x$ in $x=1$, dove la derivata esatta vale
$\cos1$ (lo dimostreremo nel \cref{sec:derivate-elementari}), confrontando la
\emph{differenza in avanti} $\frac{f(x+h)-f(x)}{h}$ con la \emph{differenza centrata}
$\frac{f(x+h)-f(x-h)}{2h}$, che ha anch'essa limite $f'(x)$ (\cref{es:differenza-centrata}).

\begin{python}[Derivata numerica di $\sen x$ in $x=1$]
import math

x, esatta = 1.0, math.cos(1.0)        # f = sen, f'(1) = cos 1
print("   h       errore in avanti   errore centrato")
for k in range(1, 13, 1):
    h = 10.0**(-k)
    avanti = (math.sin(x + h) - math.sin(x)) / h
    centrata = (math.sin(x + h) - math.sin(x - h)) / (2 * h)
    ea, ec = abs(avanti - esatta), abs(centrata - esatta)
    print(f"1e-{k:<2d}   {ea:14.3e}   {ec:14.3e}")
\end{python}
\begin{risultato}
   h       errore in avanti   errore centrato
1e-1         4.294e-02        9.001e-04
1e-2         4.216e-03        9.005e-06
1e-3         4.208e-04        9.005e-08
1e-4         4.207e-05        9.004e-10
1e-5         4.207e-06        1.114e-11
1e-6         4.207e-07        2.772e-11
1e-7         4.183e-08        1.943e-10
1e-8         2.970e-09        2.581e-09
1e-9         5.254e-08        2.970e-09
1e-10        5.848e-08        5.848e-08
1e-11        1.169e-06        1.169e-06
1e-12        4.324e-05        1.227e-05
\end{risultato}

\begin{idea}[Troncamento contro arrotondamento: le differenze finite]
L'errore della derivata numerica ha due componenti con andamento opposto: al diminuire
di $h$ la prima decresce, la seconda cresce (\cref{fig:derivata-numerica}).
L'\emph{errore di troncamento} è dovuto al fatto che $h$ non è nullo: per la differenza in avanti è proporzionale ad $h$ (lo si vede
nell'output: dividendo $h$ per $10$ l'errore si divide per $10$), per la differenza
centrata è proporzionale ad $h^{2}$ (l'errore si divide per $100$), come si dimostra con
la formula di Taylor (\cref{es:errore-differenze}). L'\emph{errore di arrotondamento} è dovuto al fatto che il
calcolatore rappresenta $f(x+h)$ e $f(x)$ con circa $16$ cifre significative: la loro
differenza perde tante più cifre quanto più i due numeri sono vicini (è la
cancellazione del \cref{cap:funzioni-continue}), e l'errore assoluto, dell'ordine di
$10^{-16}$, viene poi diviso per $h$. Il risultato è che esiste un passo ottimale, circa
$10^{-8}$ per la differenza in avanti e circa $10^{-5}$ per quella centrata, sotto il
quale \emph{diminuire $h$ peggiora il risultato}. È esattamente il dilemma di chi scrive
un codice alle differenze finite per la fluidodinamica computazionale: raffinare la
griglia riduce l'errore di discretizzazione, ma non all'infinito.
\end{idea}

\begin{figure}[htb]
\centering
\includegraphics[width=0.92\textwidth]{derivata_numerica}
\caption{Errore della derivata numerica di $\sen x$ in $x=1$ in funzione del passo $h$,
in scala bilogaritmica: a destra domina l'errore di troncamento (rette di pendenza $1$ e
$2$), a sinistra quello di arrotondamento, che cresce come $\varepsilon/h$ con
$\varepsilon\approx10^{-16}$ (\texttt{python/fig\_derivata\_numerica.py}).}
\label{fig:derivata-numerica}
\end{figure}

\subsection{Derivate successive}

Se una funzione è derivabile in tutti i punti di un intervallo $(a,b)$, allora la sua
derivata $f'(x)$ è una funzione definita su $(a,b)$. Se questa funzione è a sua volta
derivabile, diremo che la sua derivata $(f')'$ è la \emph{derivata seconda} della
funzione $f$, e indicheremo tale derivata con uno dei simboli
\begin{equation}\label{eq:notazioni-derivata-seconda}
  f'', \qquad \frac{d^{2}f}{dx^{2}}, \qquad D^{2}f, \qquad y'', \qquad \frac{d^{2}y}{dx^{2}},
  \qquad D^{2}y .
\end{equation}
Se a sua volta la derivata seconda è derivabile parleremo di derivata terza $f'''$, e
così via. Useremo il simbolo $f^{(n)}$ per la derivata $n$-esima (con la convenzione
$f^{(0)}=f$). Può accadere che una funzione ammetta derivate fino a un ordine
$n\in\N$; oppure che sia derivabile infinite volte, cioè che ammetta derivate di
qualsiasi ordine. Se la derivata $n$-esima esiste ed è continua in un intervallo, si
dice che $f$ è \emph{di classe $C^{n}$} in quell'intervallo; se $f$ ha derivate di
ogni ordine si dice che è \emph{di classe $C^{\infty}$}. Le funzioni continue si dicono
anche di classe $C^{0}$.

\begin{esempio}\label{es:x-modulo-x}
La funzione $f(x)=x^{2}$ ammette derivate di ogni ordine; infatti abbiamo verificato
che $f'(x)=2x$; poi dalla \eqref{eq:derivata-retta} segue che $f''(x)=2$ e dalla
\eqref{eq:derivata-costante} che $f'''(x)=0$; dato che $f'''$ è costante, si ottiene
$f^{(4)}=0$ e, analogamente, $f^{(n)}(x)=0$ per ogni $n\ge3$.

Invece la funzione $f(x)=x\abs x$ ammette per $x=0$ derivata prima, ma non derivata
seconda. Infatti, essendo
\begin{equation}\label{eq:x-modulo-x}
  f(x)=
  \begin{cases}
    x^{2} & \text{se } x\ge0,\\
    -x^{2} & \text{se } x<0,
  \end{cases}
\end{equation}
$f$ coincide localmente con $x^{2}$ o con $-x^{2}$ in un intorno di ogni $x\neq0$, ed è
quindi derivabile per $x\neq0$, con
\begin{equation}\label{eq:x-modulo-x-derivata}
  f'(x)=
  \begin{cases}
    2x & \text{se } x>0,\\
    -2x & \text{se } x<0,
  \end{cases}
\end{equation}
mentre, se $x=0$, si ha
\begin{equation}\label{eq:x-modulo-x-zero}
  \lim_{h\to0}\frac{f(0+h)-f(0)}{h}=\lim_{h\to0}\frac{h\abs h-0}{h}=\lim_{h\to0}\abs h=0;
\end{equation}
pertanto $f'(0)=0$. In definitiva $f$ è derivabile per ogni $x\in\R$ e la derivata vale
\begin{equation}\label{eq:x-modulo-x-finale}
  f'(x)=2\abs x \qquad\perogni x\in\R .
\end{equation}
Però $f'(x)=2\abs x$ non è derivabile per $x=0$ (\cref{es:modulo-non-derivabile});
quindi non esiste la derivata seconda di $f$ nel punto $x=0$. La funzione $x\abs x$ è di
classe $C^{1}$, ma non è derivabile due volte su $\R$.
\end{esempio}

\needspace{10\baselineskip}
\begin{notastorica}[Da Lagrange a Cauchy: il nome e la definizione]
Né Newton né Leibniz disponevano di una nozione chiara di limite, e il loro modo di
procedere, dividere per un incremento e poi porlo uguale a zero, fu duramente criticato
(si veda la nota storica del \S\ref{sec:def-limite}). Joseph-Louis Lagrange (1736--1813)
tentò di aggirare la difficoltà con l'algebra. Nella \emph{Théorie des fonctions
analytiques} (1797) supponeva che ogni funzione si potesse sviluppare in serie di potenze
dell'incremento, $f(x+i)=f(x)+p\,i+q\,i^{2}+r\,i^{3}+\cdots$, e chiamava il coefficiente
$p$, che è a sua volta una funzione di $x$, \emph{fonction dérivée} di $f$, scrivendolo
$f'x$; ripetendo l'operazione otteneva $f''x$, $f'''x$, \ldots. Da qui vengono il nome
«derivata» e gli apici delle derivate successive. Ma l'esistenza dello sviluppo non si
può giustificare senza sapere già che cosa sia una derivata: il programma conteneva un
circolo vizioso, e ne riparleremo nel \cref{cap:taylor}. La via giusta era quella
indicata da d'Alembert nell'\emph{Encyclopédie}: fondare il calcolo sul limite.
Augustin-Louis Cauchy la percorse fino in fondo nel \emph{Résumé des leçons sur le calcul
infinitésimal} (1823), le lezioni dell'École Polytechnique: la derivata vi è definita, come
nella \cref{def:derivata}, come limite del rapporto incrementale, e di Lagrange restano il
nome e il simbolo. Si capovolge anche l'ordine tra derivata e differenziale: Leibniz
definiva il differenziale per mezzo della tangente e lo trattava di fatto come un
incremento infinitesimo; per Cauchy l'oggetto primitivo è la derivata, $dx$ è una
quantità finita arbitraria e il differenziale è $dy=f'(x)\,dx$, come nel
\cref{sec:retta-tangente}.
\end{notastorica}
